\documentclass[10pt,a4paper]{amsart}
\usepackage{amsmath,amssymb,amsthm,graphicx}
\usepackage{verbatim}
\newenvironment{reference}{\comment}{\endcomment}
%\newenvironment{reference}{}{}
\newenvironment{slogan}{\comment}{\endcomment}
\newenvironment{history}{\comment}{\endcomment}

% For commutative diagrams we use Xy-pic
\usepackage[all]{xy}

% We use 2cell for 2-commutative diagrams.
\xyoption{2cell}
% Theorem environments.
%
\theoremstyle{plain}
\newtheorem{theorem}[subsection]{Theorem}
\newtheorem{proposition}[subsection]{Proposition}
\newtheorem{lemma}[subsection]{Lemma}

\theoremstyle{definition}
\newtheorem{definition}[subsection]{Definition}
\newtheorem{example}[subsection]{Example}
\newtheorem{exercise}[subsection]{Exercise}
\newtheorem{situation}[subsection]{Situation}

\theoremstyle{remark}
\newtheorem{remark}[subsection]{Remark}
\newtheorem{remarks}[subsection]{Remarks}

\numberwithin{equation}{subsection}

% Macros
%
\def\lim{\mathop{\mathrm{lim}}\nolimits}
\def\colim{\mathop{\mathrm{colim}}\nolimits}
\def\Spec{\mathop{\mathrm{Spec}}}
\def\Hom{\mathop{\mathrm{Hom}}\nolimits}
\def\Ext{\mathop{\mathrm{Ext}}\nolimits}
\def\SheafHom{\mathop{\mathcal{H}\!\mathit{om}}\nolimits}
\def\SheafExt{\mathop{\mathcal{E}\!\mathit{xt}}\nolimits}
\def\Sch{\mathit{Sch}}
\def\Mor{\mathop{\mathrm{Mor}}\nolimits}
\def\Ob{\mathop{\mathrm{Ob}}\nolimits}
\def\Sh{\mathop{\mathit{Sh}}\nolimits}
\def\NL{\mathop{N\!L}\nolimits}
\def\CH{\mathop{\mathrm{CH}}\nolimits}
\def\proetale{{pro\text{-}\acute{e}tale}}
\def\etale{{\acute{e}tale}}
\def\QCoh{\mathit{QCoh}}
\def\Ker{\mathop{\mathrm{Ker}}}
\def\Im{\mathop{\mathrm{Im}}}
\def\Coker{\mathop{\mathrm{Coker}}}
\def\Coim{\mathop{\mathrm{Coim}}}

% Boxtimes
%
\DeclareMathSymbol{\boxtimes}{\mathbin}{AMSa}{"02}

%
% Macros for moduli stacks/spaces
%
\def\QCohstack{\mathcal{QC}\!\mathit{oh}}
\def\Cohstack{\mathcal{C}\!\mathit{oh}}
\def\Spacesstack{\mathcal{S}\!\mathit{paces}}
\def\Quotfunctor{\mathrm{Quot}}
\def\Hilbfunctor{\mathrm{Hilb}}
\def\Curvesstack{\mathcal{C}\!\mathit{urves}}
\def\Polarizedstack{\mathcal{P}\!\mathit{olarized}}
\def\Complexesstack{\mathcal{C}\!\mathit{omplexes}}
% \Pic is the operator that assigns to X its picard group, usage \Pic(X)
% \Picardstack_{X/B} denotes the Picard stack of X over B
% \Picardfunctor_{X/B} denotes the Picard functor of X over B
\def\Pic{\mathop{\mathrm{Pic}}\nolimits}
\def\Picardstack{\mathcal{P}\!\mathit{ic}}
\def\Picardfunctor{\mathrm{Pic}}
\def\Deformationcategory{\mathcal{D}\!\mathit{ef}}
\usepackage{cleveref}
\crefname{originaldefinition}{original definition}{original definitions}
\Crefname{originaldefinition}{Original definition}{Original definitions}
\makeatletter
\newlabel{definition-etale-smooth-local-source-target}{{0CFZ}{source}{}{sourcepointer-1}{}}
\newlabel{definition-etale-smooth-local-source-target@cref}{{[originaldefinition][0][]0CFZ}{[1][0][]source}{}{}{}}
\makeatother
\crefname{originaldefinition}{original definition}{original definitions}
\Crefname{originaldefinition}{Original definition}{Original definitions}
\makeatletter
\newlabel{definition-local-source-target}{{06F8}{source}{}{sourcepointer-2}{}}
\newlabel{definition-local-source-target@cref}{{[originaldefinition][0][]06F8}{[1][0][]source}{}{}{}}
\makeatother
\crefname{originalreference}{original reference}{original references}
\Crefname{originalreference}{Original reference}{Original references}
\makeatletter
\newlabel{spaces-properties-lemma-etale-permanence}{{03FV}{source}{}{sourcepointer-3}{}}
\newlabel{spaces-properties-lemma-etale-permanence@cref}{{[originalreference][0][]03FV}{[1][0][]source}{}{}{}}
\makeatother
\crefname{originalreference}{original reference}{original references}
\Crefname{originalreference}{Original reference}{Original references}
\makeatletter
\newlabel{spaces-morphisms-lemma-etale-permanence}{{05W3}{source}{}{sourcepointer-4}{}}
\newlabel{spaces-morphisms-lemma-etale-permanence@cref}{{[originalreference][0][]05W3}{[1][0][]source}{}{}{}}
\makeatother
\title{Permanence properties of \textbackslash{}'etale-smooth local on source-and-target properties: \textbackslash{}'etale locality on the source, smooth locality on the target, and postcomposition/permanence for \textbackslash{}'etale morphisms}
\author{The Stacks Project Authors}
\date{stacks-project:spaces-descent.tex:local-snapshot-20261012; GNU FDL 1.2 or later}
\begin{document}
\maketitle
\noindent\textit{Editorial note.} All delivered mathematical text is the author's own, copied verbatim from the pinned Stacks Project chapter file spaces-descent.tex; no statement or proof step was rewritten, reordered or removed.\par\medskip
\noindent\textit{Editorial note.} The author's opening sentence 'We write everything out completely.' is part of the counted proof text and is kept verbatim.\par\medskip
\noindent\textit{Editorial note.} The counted span cites two definitions and two permanence lemmas, recorded in assembly.external\_references with canonical tags as navigation pointers.\par\medskip
\noindent\textit{Allowed standard prerequisites (editorial).}\begin{itemize}
\item \textbackslash{}'Etale-smooth local on source-and-target: definition. Stacks Project, Descent on Spaces, Definition definition-etale-smooth-local-source-target. Supplies parts (1), (2), (3) used to unfold and re-fold property P along coverings.
\item Smooth local on source-and-target: definition. Stacks Project, Descent on Spaces, Definition definition-local-source-target. Supplies the target-locality unfolding used in parts (2) and (3).
\item \textbackslash{}'Etale permanence for algebraic spaces. Stacks Project, Properties of Spaces, Lemma spaces-properties-lemma-etale-permanence; Morphisms of Spaces, Lemma spaces-morphisms-lemma-etale-permanence. Supplies that X \textbackslash{}to Y \textbackslash{}times\_Z X is \textbackslash{}'etale in part (3) and that the section (f, 1) is \textbackslash{}'etale in part (4).
\end{itemize}\medskip
\ifdefined\hypertarget\hypertarget{sourcepointer-1}{}\fi
\noindent\textit{Source reference.} Original definition 0CFZ: sources/stacks-chow/tags.txt, source lines 15629--15629. Cross-chapter standard prerequisite cited by the author inside the counted span; it is used as a pointer and its own proof is not part of this item.\par\medskip
\ifdefined\hypertarget\hypertarget{sourcepointer-2}{}\fi
\noindent\textit{Source reference.} Original definition 06F8: sources/stacks-chow/tags.txt, source lines 8242--8242. Cross-chapter standard prerequisite cited by the author inside the counted span; it is used as a pointer and its own proof is not part of this item.\par\medskip
\ifdefined\hypertarget\hypertarget{sourcepointer-3}{}\fi
\noindent\textit{Source reference.} Original reference 03FV: sources/stacks-chow/tags.txt, source lines 4506--4506. Cross-chapter standard prerequisite cited by the author inside the counted span; it is used as a pointer and its own proof is not part of this item.\par\medskip
\ifdefined\hypertarget\hypertarget{sourcepointer-4}{}\fi
\noindent\textit{Source reference.} Original reference 05W3: sources/stacks-chow/tags.txt, source lines 7572--7572. Cross-chapter standard prerequisite cited by the author inside the counted span; it is used as a pointer and its own proof is not part of this item.\par\medskip
\begin{lemma}
\label{lemma-etale-smooth-local-source-target-implies}
Let $S$ be a scheme.
Let $\mathcal{P}$ be a property of morphisms of algebraic spaces over $S$
which is \'etale-smooth local on source-and-target. Then
\begin{enumerate}
\item $\mathcal{P}$ is \'etale local on the source,
\item $\mathcal{P}$ is smooth local on the target,
\item $\mathcal{P}$ is stable under postcomposing with \'etale morphisms:
if $f : X \to Y$ has $\mathcal{P}$ and $g : Y \to Z$ is \'etale, then
$g \circ f$ has $\mathcal{P}$, and
\item $\mathcal{P}$ has a permanence property: given $f : X \to Y$ and
$g : Y \to Z$ \'etale such that $g \circ f$ has $\mathcal{P}$, then
$f$ has $\mathcal{P}$.
\end{enumerate}
\end{lemma}

\begin{proof}
We write everything out completely.

\medskip\noindent
Proof of (1). Let $f : X \to Y$ be a morphism of algebraic spaces over $S$.
Let $\{X_i \to X\}_{i \in I}$ be an \'etale covering of $X$. If each composition
$h_i : X_i \to Y$ has $\mathcal{P}$, then for each $|x| \in X$ we can find
an $i \in I$ and a point $x_i \in |X_i|$ mapping to $x$. Then
$(X_i, x_i) \to (X, x)$ is an \'etale morphism of pairs, and
$\text{id}_Y : Y \to Y$ is a smooth morphism, and $h_i$ is as in part (3) of
Definition \ref{definition-etale-smooth-local-source-target}.
Thus we see that $f$ has $\mathcal{P}$.
Conversely, if $f$ has $\mathcal{P}$ then each $X_i \to Y$ has
$\mathcal{P}$ by
Definition \ref{definition-etale-smooth-local-source-target} part (1).

\medskip\noindent
Proof of (2). Let $f : X \to Y$ be a morphism of algebraic spaces over $S$.
Let $\{Y_i \to Y\}_{i \in I}$ be a smooth covering of $Y$.
Write $X_i = Y_i \times_Y X$ and $h_i : X_i \to Y_i$ for the base change
of $f$. If each  $h_i : X_i \to Y_i$ has $\mathcal{P}$, then for each
$x \in |X|$ we pick an $i \in I$ and a point $x_i \in |X_i|$ mapping to $x$.
Then $X_i \to X \times_Y Y_i$ is an \'etale morphism
(because it is an isomorphism), $Y_i \to Y$ is
smooth, and $h_i$ is as in part (3) of
Definition \ref{definition-local-source-target}.
Thus we see that $f$ has $\mathcal{P}$.
Conversely, if $f$ has $\mathcal{P}$, then each $X_i \to Y_i$ has
$\mathcal{P}$ by
Definition \ref{definition-local-source-target} part (2).

\medskip\noindent
Proof of (3). Assume $f : X \to Y$ has $\mathcal{P}$ and $g : Y \to Z$ is
\'etale. The morphism $X \to Y \times_Z X$ is \'etale as a morphism
between algebraic spaces \'etale over $X$ (
Properties of Spaces, Lemma \ref{spaces-properties-lemma-etale-permanence}).
Also $Y \to Z$ is \'etale hence a smooth morphism.
Thus the diagram
$$
\xymatrix{
X \ar[d] \ar[r]_f & Y \ar[d] \\
X \ar[r]^{g \circ f} & Z
}
$$
works for every $x \in |X|$ in part (3) of
Definition \ref{definition-local-source-target}
and we conclude that $g \circ f$ has $\mathcal{P}$.

\medskip\noindent
Proof of (4). Let $f : X \to Y$ be a morphism and $g : Y \to Z$ \'etale
such that $g \circ f$ has $\mathcal{P}$. Then by
Definition \ref{definition-etale-smooth-local-source-target} part (2)
we see that $\text{pr}_Y : Y \times_Z X \to Y$ has $\mathcal{P}$. But
the morphism $(f, 1) : X \to Y \times_Z X$ is \'etale as a section to the
\'etale projection $\text{pr}_X : Y \times_Z X \to X$, see
Morphisms of Spaces, Lemma \ref{spaces-morphisms-lemma-etale-permanence}.
Hence $f = \text{pr}_Y \circ (f, 1)$ has $\mathcal{P}$ by
Definition \ref{definition-etale-smooth-local-source-target} part (1).
\end{proof}

\end{document}
